\documentclass[10pt]{article}
\providecommand{\FIGDIR}{fig}
\newcommand{\DIGESTDATE}{16 September 2026}
% ===================== PM Research Watch - shared preamble =====================
\usepackage[T1]{fontenc}
\usepackage[utf8]{inputenc}
\usepackage{amsmath,amssymb}
\usepackage{tgpagella}
\usepackage{tgheros}
\renewcommand{\sfdefault}{qhv}
\renewcommand{\ttdefault}{lmtt}
\usepackage{microtype}
\usepackage[english]{babel}

\usepackage[a4paper,top=17mm,bottom=17mm,left=17mm,right=17mm,
            headsep=6mm,footskip=10mm]{geometry}
\usepackage{graphicx}
\usepackage{xcolor}
\usepackage{colortbl}
\usepackage{tikz}
\usetikzlibrary{calc}
\usepackage[most]{tcolorbox}
\usepackage{titlesec}
\usepackage{fancyhdr}
\usepackage{booktabs}
\usepackage{array}
\usepackage{longtable}
\usepackage{enumitem}
\usepackage{multicol}
\usepackage{ragged2e}
\usepackage[hidelinks]{hyperref}

% ------------------------------------------------------------------ palette
\definecolor{Ink}    {HTML}{14262E}
\definecolor{Deep}   {HTML}{0F4C5C}
\definecolor{Teal}   {HTML}{1F7A8C}
\definecolor{Sky}    {HTML}{4FA3C4}
\definecolor{Sage}   {HTML}{6FA287}
\definecolor{Amber}  {HTML}{E3A008}
\definecolor{Clay}   {HTML}{B5643C}
\definecolor{Coral}  {HTML}{D1495B}
\definecolor{Violet} {HTML}{6A5B9E}
\definecolor{Slate}  {HTML}{4A6572}
\definecolor{Paper}  {HTML}{FBF9F5}
\definecolor{Mist}   {HTML}{EEEAE2}
\definecolor{Rule}   {HTML}{D9D3C9}

\pagecolor{Paper}
\color{Ink}

\hypersetup{colorlinks=true, linkcolor=Deep, urlcolor=Teal, citecolor=Teal,
            pdftitle={PM Research Watch}, pdfauthor={Automated daily digest}}

% ------------------------------------------------------------------ headings
\titleformat{\section}
  {\sffamily\bfseries\large\color{Deep}}{}{0pt}
  {}[{\vspace{1.2mm}\color{Amber}\titlerule[1.1pt]}]
\titlespacing*{\section}{0pt}{5.5mm}{2.8mm}

\titleformat{\subsection}
  {\sffamily\bfseries\normalsize\color{Teal}}{}{0pt}{}
\titlespacing*{\subsection}{0pt}{3.6mm}{1.4mm}

\setlist[itemize]{leftmargin=4.2mm, itemsep=1.1mm, topsep=1.3mm, parsep=0pt,
                  label=\textcolor{Teal}{\footnotesize$\blacktriangleright$}}

% ------------------------------------------------------------------ header/footer
\pagestyle{fancy}
\fancyhf{}
\renewcommand{\headrulewidth}{0pt}
\renewcommand{\footrulewidth}{0pt}
\fancyhead[L]{\sffamily\scriptsize\color{Slate}PM RESEARCH WATCH\;\textcolor{Rule}{\textbar}\;\DIGESTDATE}
\fancyhead[R]{\sffamily\scriptsize\color{Slate}Particulate matter monitoring \& health effects}
\fancyfoot[C]{\sffamily\scriptsize\color{Slate}\thepage}

% ------------------------------------------------------------------ components
% metric tile
\newcommand{\metric}[3]{%
  \begin{tcolorbox}[enhanced, colback=white, colframe=Rule, boxrule=0.5pt,
    arc=1.6mm, left=2.2mm, right=2.2mm, top=1.8mm, bottom=1.8mm,
    borderline west={1.7mm}{0pt}{#3}, nobeforeafter, valign=center,
    height=17mm, width=\linewidth]
    \RaggedRight
    {\sffamily\bfseries\fontsize{17}{17}\selectfont\color{#3}#1}\\[0.5mm]
    {\sffamily\fontsize{6.5}{7.4}\selectfont\color{Slate}#2}
  \end{tcolorbox}}

% pull-quote / key insight box
\newtcolorbox{keybox}[1][]{enhanced, colback=Mist, colframe=Mist,
  arc=1.6mm, left=3.2mm, right=3.2mm, top=2.4mm, bottom=2.4mm,
  borderline west={1.7mm}{0pt}{Amber}, #1}

% subtopic band
\newcommand{\band}[2]{%
  \vspace{3.2mm}\noindent
  \begin{tcolorbox}[enhanced, colback=#2, colframe=#2, arc=1.2mm,
    left=2.6mm, right=2.6mm, top=1.1mm, bottom=1.1mm, boxrule=0pt,
    width=\linewidth, nobeforeafter]
    {\sffamily\bfseries\small\color{white}#1}
  \end{tcolorbox}\vspace{1.6mm}}

% one paper entry: {short cite}{title}{journal}{doi}{body}
\newcommand{\paperentry}[5]{%
  \par\vspace{1.9mm}\noindent
  {\sffamily\bfseries\footnotesize\color{Ink}#1}\;%
  {\sffamily\fontsize{7}{8}\selectfont\color{Slate}#3}\\[0.4mm]
  {\footnotesize\itshape\color{Deep}#2}\\[0.6mm]
  {\footnotesize #4}\\[0.3mm]
  {\fontsize{6.6}{7.6}\selectfont\color{Slate}\texttt{doi:}\,\href{https://doi.org/#5}{\texttt{\detokenize{#5}}}}
  \par}

% caption style
\newcommand{\figcap}[1]{\par\vspace{1.0mm}
  {\fontsize{7.2}{8.6}\selectfont\color{Slate}\textsf{\textbf{Fig.}}\; #1\par}\vspace{1.5mm}}

\providecommand{\permil}{\textperthousand}

\setlength{\parindent}{0pt}
\setlength{\parskip}{0pt}
\renewcommand{\arraystretch}{1.18}


\begin{document}

% ==================================================================== MASTHEAD
\thispagestyle{empty}
\begin{tikzpicture}[remember picture, overlay]
  \fill[Deep] (current page.north west) rectangle
      ($(current page.north east)+(0,-67mm)$);
  \fill[Amber] ($(current page.north west)+(0,-67mm)$) rectangle
      ($(current page.north east)+(0,-68.4mm)$);
  \foreach \x/\y/\r in {22/12/0.9, 41/26/1.5, 63/9/0.7, 88/31/1.1, 112/17/1.9,
                        138/29/0.8, 159/13/1.3, 178/54/1.0, 196/21/1.6,
                        30/44/0.6, 74/57/1.2, 125/48/0.9, 168/6/0.8, 191/59/0.7,
                        52/18/1.0, 100/60/0.6, 148/20/1.1, 66/49/1.4, 205/45/0.8}
    \fill[Sky, opacity=0.30] ($(current page.north west)+(\x mm,-\y mm)$) circle (\r mm);
\end{tikzpicture}

\vspace*{4mm}
{\sffamily\fontsize{7.6}{9}\selectfont\color{Sky}%
 AUTOMATED DAILY LITERATURE DIGEST\;\textbullet\;PhD RESEARCH WATCH\par}
\vspace{2.2mm}
{\sffamily\bfseries\fontsize{27}{29}\selectfont\color{white}%
 Particulate Matter\par}
{\sffamily\fontsize{27}{29}\selectfont\color{Sky}%
 Monitoring \& Health Effects\par}
\vspace{5.0mm}
{\sffamily\fontsize{8.4}{10}\selectfont\color{white}%
 Issue 2026-09-16\;\;\textcolor{Amber}{\textbar}\;\;PubMed entry window 16 September 2026\;\;%
 \textcolor{Amber}{\textbar}\;\;21 records screened in scope\par}
\vspace{18mm}

% -------------------------------------------------------------- metric strip
\noindent
\begin{minipage}[t]{0.190\linewidth}\metric{21}{RECORDS IN SCOPE\\(27 EXCLUDED, 3 HELD)}{Deep}\end{minipage}\hfill
\begin{minipage}[t]{0.190\linewidth}\metric{9}{OF 10 SUBTOPIC\\BANDS POPULATED}{Teal}\end{minipage}\hfill
\begin{minipage}[t]{0.190\linewidth}\metric{6}{EFFECT ESTIMATES\\(ALL WITH A CI)}{Sky}\end{minipage}\hfill
\begin{minipage}[t]{0.190\linewidth}\metric{5}{TIER-A\\RECORDS}{Amber}\end{minipage}\hfill
\begin{minipage}[t]{0.190\linewidth}\metric{0}{CONSENSUS RECORDS\\ADMITTED}{Clay}\end{minipage}

\vspace{6mm}

% -------------------------------------------------------------- signal of day
\section*{\sffamily\bfseries\color{Deep}Signal of the day}
\vspace{-1.5mm}{\color{Amber}\rule{\linewidth}{1.1pt}}\vspace{2.5mm}

\begin{keybox}
{\sffamily\bfseries\small\color{Deep}Landscape-fire PM$_{2.5}$ is 19\% of ambient mass
and 31\% of the attributable deaths --- and conventional epidemiology has been
understating it by more than half.}\\[1.4mm]
\footnotesize
Ju et al.\ (\textit{Lancet Planet Health}) linked daily mortality in \textbf{1{,}366
communities across ten countries and territories}, 2000--19, \textbf{47{,}844{,}851
deaths}, to population-weighted landscape-fire PM$_{2.5}$, and estimated the association
with instrumental-variable models rather than the usual time-series regression. Each
1\,$\mu$g/m$^3$ of 2-day mean fire PM$_{2.5}$ raised mortality
\textbf{0.60\% (95\% CI 0.19--1.00)} and mortality-related economic losses
\textbf{0.59\% (0.21--0.97)}, giving \textbf{906{,}869 attributable deaths} and
\textbf{US\$994.5 billion}. Two numbers carry the issue. First, the disproportion:
fire-sourced particles are \textbf{19.0\% of ambient PM$_{2.5}$ mass but 30.6\% of
PM$_{2.5}$-attributable deaths} --- roughly a \textbf{1.9-fold per-unit potency premium}
over background mass. Second, the method gap: the IV estimates exceeded conventional
association estimates by a \textbf{median 56.0\% (2.2--111.3)} for mortality and
\textbf{62.3\% (9.0--116.1)} for economic loss, i.e.\ the standard design is biased
\emph{toward the null} here, not away from it. The instrument itself is not named in the
abstract, and exclusion-restriction validity against co-emitted gases and weather carries
the whole causal claim, so this is a strong result resting on one unexamined assumption.
For instrumentation the implication is blunt: a network reporting total PM$_{2.5}$
mass cannot see the fraction that carries a third of the deaths.
\end{keybox}

\clearpage

% -------------------------------------------------------------- exec brief
\section*{\sffamily\bfseries\color{Deep}What changed today}
\vspace{-1.5mm}{\color{Amber}\rule{\linewidth}{1.1pt}}\vspace{2.5mm}

\footnotesize
\begin{itemize}
\item \textbf{Two independent causal-inference papers both found the tail, not the mean,
carries the burden.} Ju et al.\ show fire-sourced PM$_{2.5}$ is potency-enriched relative
to its mass share. Bao et al.\ (\textit{J Environ Manage}), evaluating 88 Chinese
clean-heating pilot cities with an augmented synthetic control, found the policy moved the
heating-season \emph{mean} only $-$3.01\,$\mu$g/m$^3$ but P95 by $-$5.76, P99 by $-$5.96
and exceedance days by $-$7.92; high-pollution days were 27.8\% of improvement days and
48.8\% of the mortality difference. A mean-only evaluation of either paper's exposure
would have halved its own answer. Both argue for networks resolved enough to capture
episode tails.

\item \textbf{Aerosol-layer height and fluorescence lidar both landed in the same window,
and the second undermines a standing sensor assumption.} Li et al.\ (\textit{AMT})
retrieve aerosol layer height from HARP2/PACE multi-angle polarimetry at
\textbf{RMSE 1.03\,km} against EarthCARE ATLID --- the vertical term that dominates
AOD-to-surface-PM error for smoke and dust. Hu et al.\ (\textit{ACP}), observing 2023
Canadian wildfire smoke at two sites with different fluorescence lidar designs, report
\textbf{near-absence of hygroscopic growth at 90--100\% RH}. If that generalises, every
$\kappa$-K\"ohler humidity correction applied to optical low-cost PM sensors is
systematically wrong during exactly the episodes those sensors are deployed to catch.

\item \textbf{A clean PM$_{2.5}$ null, which the archive needs more of.} Ammons et al.\
(\textit{Environ Res}) meta-analysed seven studies of PM$_{2.5}$ and gynaecologic cancer:
ovarian \textbf{HR 1.02 (0.91--1.13)} per 5\,$\mu$g/m$^3$, and \textbf{0.94 (0.81--1.09)}
restricted to co-pollutant-adjusted studies. The endometrial estimate is nominally
positive but its confidence interval --- 1.01 (1.00--1.01) --- is implausibly tight for
five studies and no $I^2$ is reported, which points to shared cohorts rather than
precision.

\item \textbf{Consensus contributed nothing admissible for the first time.} The sweep
returned 10 calibration records and 10 fusion records; \textbf{seven were already in the
archive} by title, and \textbf{every one of the five that survived title screening failed a
Crossref \texttt{created} check} --- 2026-01-02, 2026-05-08, 2026-05-27, 2026-08-22 and
2025-12-23. See the provenance note.

\item \textbf{Europe PMC's AGRICOLA leg was found to retro-index.} Two records returned
inside a \texttt{CREATION\_DATE} window of 16 September resolve to \textbf{pubYear 2021 and
2014}. This is a source defect, not a screening miss, and it is the reason the AGR source
now requires a publication-year check before any record from it is summarised.

\item \textbf{Indoor and appliance-side exposure got two usable methods papers.} Tam et
al.\ (\textit{Curr Environ Health Rep}) find nonzero emissions for \textbf{at least 41
toxic air contaminants} from gas appliances, with formaldehyde and benzene emission
factors spanning \textbf{two orders of magnitude} --- and suggest CH$_4$ and CO as
inference surrogates, which is a directly implementable low-cost sensing target.
Yoshizaki et al.\ (\textit{J Am Coll Surg}), in a 64-patient randomised trial, cut
particle counts \textbf{96--97\%} with an 8\,mm capture distance and found no significant
effect at 18\,mm --- a 10\,mm change in near-field geometry deciding the whole result.
\end{itemize}

\clearpage

% ==================================================================== ANALYTICS
\section*{\sffamily\bfseries\color{Deep}Corpus analytics}
\vspace{-1.5mm}{\color{Amber}\rule{\linewidth}{1.1pt}}\vspace{3mm}

\noindent\includegraphics[width=\linewidth]{\FIGDIR/f1_subtopics.png}
\figcap{Subtopic assignment of the 21 in-scope records; single-label by dominant
contribution. \textit{Sensing, forecasting \& instrumentation} leads with 5 and
\textit{Exposure assessment \& modelling} follows with 4, so the two measurement-side
bands together hold 9 of 21 --- the instrumentation axis outweighs any single health axis
today. \textit{Burden, policy \& mitigation} holds 3, three bands hold 2 (respiratory,
occupational/indoor, mechanistic) and three hold a single record each.
\textit{Other clinical endpoints} is the one empty band.}

\vspace{2mm}
\noindent\includegraphics[width=\linewidth]{\FIGDIR/f2_design_endpoint.png}
\figcap{Left --- study architecture. \textit{Review / synthesis} is the largest group at
6 of 21 (29\%), an unusually review-heavy day; \textit{Modelling / inventory} holds 4,
\textit{Experimental / toxicology} and \textit{Measurement campaign} 3 each,
\textit{Ecological} and \textit{Observational --- cross-sectional} 2 each, and there is a
single randomised trial. Right --- health endpoint. \textbf{11 of 21 records carry no
human health endpoint at all}, the single largest category and a direct consequence of the
measurement-side weighting above; \textit{Respiratory} is the largest health endpoint at 4,
and mortality, cardiovascular, neuro/mental-health, oncologic and in-vitro toxicology hold
one or two each.}

\vspace{2mm}
\noindent\includegraphics[width=\linewidth]{\FIGDIR/f3_metric_geography.png}
\figcap{Left --- particle metric. Unspeciated PM and emissions lead at 7, PM$_{2.5}$ alone
at 6, and \textbf{optical properties (absorption/scattering) at 4} --- the AMT and ACP
retrieval records, an unusually large optical share for a daily issue. Ultrafine particles
are the explicit subject of 2, bioaerosol and composition 1 each. Right --- geographic
provenance. Global/multi-region leads with 9, China follows with 6, South Asia holds 2,
and Europe, North America, the Middle East \& N. Africa and East Asia (ex-China) hold one
each. Sub-Saharan Africa and Latin America are absent.}

\vspace{2mm}
\noindent\includegraphics[width=\linewidth]{\FIGDIR/f6_lifecourse.png}
\figcap{Life-course windows addressed; counts are non-exclusive. \textbf{Three of the five
windows are empty today} --- no prenatal, childhood or adolescent record appears at all,
and the one asthma model in the corpus is murine rather than paediatric. The occupational
window holds 2 (hospital-staff breath biomonitoring and the surgical-smoke trial) and older
adults 1 (the CLHLS depression cohort). A developmental-exposure gap this complete is
unusual and is worth watching across the week rather than reading from one day.}

\vspace{2mm}
\noindent\includegraphics[width=\linewidth]{\FIGDIR/f4_heatmap.png}
\figcap{Subtopic $\times$ study-architecture cross-tabulation. Two features are worth
naming precisely. First, the measurement--health split is nearly, but not quite, total:
all 5 \textit{Sensing} records carry no human health endpoint, and of the 4
\textit{Exposure assessment} records \textbf{exactly one} does --- the Haifa
breath-biomonitoring study, whose endpoint is respiratory. Second,
\textit{Mechanistic toxicology} is \textbf{not} wholly experimental this issue: one of its
two records is a systematic review, so the usual all-experimental pattern is broken.
The review-heavy column otherwise sits entirely on the health-side bands.}

\vspace{2mm}
\noindent\includegraphics[width=\linewidth]{\FIGDIR/f5_forest.png}
\figcap{The 6 interval-bearing estimates in this issue, on a common logarithmic axis.
\textbf{Exposure contrasts are not mutually comparable}: the two landscape-fire estimates
are per 1\,$\mu$g/m$^3$ of a 2-day mean, converted from the published percent changes to
risk-ratio equivalents and flagged as such; the gynaecologic-cancer HRs are per
5\,$\mu$g/m$^3$ of long-term exposure; and the solid-fuel OR is a binary exposed/unexposed
contrast with no measured concentration behind it. Read direction and precision, not
relative magnitude. \textbf{Three of the six cross unity}, and all three are the
gynaecologic-cancer estimates: ovarian 1.02 (0.91--1.13), co-pollutant-adjusted ovarian
0.94 (0.81--1.09) --- the only point estimate below 1 in the issue --- and endometrial
1.02 (0.99--1.05). The three that exclude unity are the two landscape-fire estimates,
whose intervals are extremely tight because they rest on 47.8 million deaths, and the
household solid-fuel OR. The endometrial interval is so narrow that it should be read as an
artefact of shared cohorts rather than as precision.}

\clearpage

% ==================================================================== DIGEST
\section*{\sffamily\bfseries\color{Deep}Paper-level digest}
\vspace{-1.5mm}{\color{Amber}\rule{\linewidth}{1.1pt}}\vspace{1mm}

{\fontsize{7.4}{9}\selectfont\color{Slate}Ordered by cluster, largest first. Each entry
gives the methodological core, the finding that matters and the limitation that actually
threatens the inference --- not an abstract paraphrase.\par}

\band{Sensing, forecasting \& instrumentation\hfill 5 records}{Amber}

\paperentry{Li et al.\ 2026}
{A multi-angle and polarization-based retrieval algorithm for aerosol layer height of smoke and dust}
{Atmos Meas Tech}
{Optimal-estimation retrieval adding multi-angle polarimetry to conventional spectral
aerosol-layer-height methods, exploiting polarisation sensitivity to the
Rayleigh-versus-aerosol scattering contrast. Information-content analysis quantifies the
gain as increased degrees of freedom for signal over radiance-only observation --- measured,
not asserted. Applied to HARP2 on PACE and validated against ATLID on EarthCARE,
\textbf{RMSE 1.03\,km} across all collocated samples. That is the number to hold: about a
kilometre of vertical uncertainty against a 1--2\,km boundary layer constrains plume
altitude but does not resolve surface coupling. Validated against a single spaceborne lidar
with limited collocation statistics and no ground profile network.
\textbf{Why it matters here:} vertical distribution is the dominant error term in
AOD-to-surface-PM$_{2.5}$ conversion, and smoke and dust are exactly where it bites.}
{10.5194/amt-19-5843-2026}

\paperentry{Hu et al.\ 2026}
{Biomass burning aerosols during the 2023 Canadian wildfires from dual-site Raman and fluorescence lidar}
{Atmos Chem Phys}
{Dual-site observation of transported smoke, May--September 2023, at ATOLL (France,
multi-wavelength Raman with a 466\,nm fluorescence channel) and GPI (Russia, five-channel
broadband fluorescence at 355\,nm excitation). Upper-troposphere/lower-stratosphere layers
showed higher depolarisation, slightly lower lidar ratios, lower {\AA}ngstr\"om exponents
and a redshifted fluorescence peak against free-troposphere layers; fluorescence magnitude
and spectral shape agreed across two sites and two instrument designs, the strongest
evidence yet that these signatures are transferable rather than instrument-specific.
The consequential finding is negative: \textbf{near-absence of hygroscopic growth at
90--100\% RH}, contradicting the standard assumption that aged biomass-burning aerosol is
hygroscopic --- with direct consequences for every $\kappa$-K\"ohler humidity correction
applied to optical PM sensors during smoke episodes. One fire season, two sites;
altitude correlations at $r^2$ 0.61--0.68 that the authors concede proxy for injection
height and plume origin.}
{10.5194/acp-26-13027-2026}

\paperentry{Ni et al.\ 2026}
{An AI-based algorithm for retrieving aerosol optical depth and single scattering albedo from All-Sky Imager observations}
{Atmos Meas Tech}
{Machine-learning retrieval trained on co-located sun-photometer AOD and SSA at two sites,
Beijing\_PKU in East China and SGP in the central US. Independent validation gave
\textbf{$r>0.86$ for AOD at all wavelengths with RMSE below 0.25}, and for SSA
\textbf{$r=0.74$ (Beijing\_PKU) and 0.84 (SGP) with RMSE below 0.09}. Retrieving SSA from a
consumer-grade sky camera is the striking part --- absorption is normally the hardest
aerosol property to retrieve and demands a calibrated multi-channel photometer. Two sites
means two aerosol regimes and two instrument configurations, and an AOD RMSE of 0.25 is
large wherever ambient AOD is low. As a densification strategy this is the same economics
that produced PurpleAir --- cheap hardware plus a trained model substituting for a
reference instrument --- with the same unanswered transferability question.}
{10.5194/amt-19-5871-2026}

\paperentry{Zhao et al.\ 2026}
{Deep learning PM$_{2.5}$ prediction integrating improved MSTL decomposition and the FATA optimizer}
{PLoS One}
{FATA-GSPMSTL-CNN-LSTM: RFECV plus Pearson feature screening, a grid-searched
power-spectral variant of Multiple Seasonal-Trend decomposition using Loess, the Fata
Morgana metaheuristic for hyperparameter search, and a CNN-LSTM head. Evaluated separately
for heating and non-heating seasons in Guangzhou and Xianyang, and the seasonal split is
the sensible part given how differently those two cities are sourced. Beyond that it
follows a familiar pattern: four stacked components, \textbf{no ablation isolating which
one carries the gain}, and improvement measured only against author-selected baselines.
No prediction intervals, so the output is a point forecast with no stated uncertainty and
cannot drive threshold alerting. Two cities, no external-city test.}
{10.1371/journal.pone.0358271}

\paperentry{Peng et al.\ 2026}
{Inhibition of condensable particulate matter by co-firing biomass and high-sulfur coal}
{Environ Pollut}
{50\,kW one-dimensional self-sustaining furnace, high-sulfur coal co-fired with corn stover
or pine wood at 10\% and 20\% blend. Baseline condensable PM was \textbf{6.72\,mg/MJ,
83.3\% inorganic}, falling to 5.09 and 3.85\,mg/MJ for corn stover at 10\% and 20\% and to
5.36 and 4.55 for pine wood --- \textbf{a 43\% reduction at 20\% corn stover}, traced to
suppression of sulfate-associated inorganics and organic precursors. PAHs and
sulfur-containing organics fell below detection at 20\% blending.
\textbf{Why it matters here:} condensable PM is invisible to hot-filter methods and to
every optical low-cost sensor, yet can rival filterable PM in mass, so inventories built
on filterable-only sampling are biased low. One bench furnace, one coal, two feedstocks,
no atmospheric ageing --- and condensable PM measurement is method-dependent enough that
absolute values should not travel across sampling protocols.}
{10.1016/j.envpol.2026.129163}


\band{Exposure assessment \& modelling\hfill 4 records}{Amber}

\paperentry{Fireman et al.\ 2026}
{Biomonitoring of urban ultrafine particle exposure and associated health effects}
{Int J Public Health}
{99 hospital staff --- 50 in the industrial Haifa Bay area, 49 Tel Aviv controls --- with
spirometry, FeNO, exhaled breath condensate and saliva at baseline and one year. Latent
class analysis found a cluster with higher BMI, raised EBC ultrafine particle concentration
and elevated salivary LDH-calcium, predominantly Haifa residents, reporting more cough,
sputum and atopic complaints. The concept is worth tracking --- EBC particle burden as an
internal-dose biomarker would close the gap between a fixed monitor and what a person
actually inhaled --- but the execution does not yet support it. \textbf{$n=99$ with LCA on a
handful of variables will find clusters whether or not they exist}; no ambient UFP
measurements accompany the biomarker, so the external-to-internal dose relation is never
tested; and BMI loading on the same cluster makes an exposure reading hard to separate
from a metabolic one.}
{10.3389/ijph.2026.1609855}

\paperentry{Li et al.\ 2026}
{A theoretical model for predicting visibility during dust events}
{Atmosphere}
{Physically based chain from surface wind erosion through dust emission and near-surface
concentration to horizontal visibility, with an analytical concentration--visibility
relation closed heuristically and a moment-balance entrainment model carrying median
particle size and soil cohesion into the threshold friction velocity. Against three
published datasets the relation explained \textbf{$R^2 = 0.62$} and had the highest $R^2$
of the formulations compared under identical calibration. The appeal is economic:
visibility is observed almost everywhere dust matters and PM$_{10}$ usually is not, so a
defensible physical inversion turns an existing network into a concentration proxy.
$R^2=0.62$ supports event detection and screening, not exposure assignment. The authors
are explicit that \textbf{only the concentration--visibility link was evaluated} --- the
emission and entrainment stages are untested, and closing the chain needs synchronised
observations that do not currently exist.}
{10.3390/atmos17090902}

\paperentry{Paygude et al.\ 2026}
{Heavy metal exposure and source profiling from leaf-deposited particulate matter in Pune}
{Environ Monit Assess}
{27 roadside, industrial, peri-urban, highway and rural sites. Elemental order
Ca $>$ K $>$ Fe $>$ Na $>$ Mg $>$ Mn $>$ Zn $>$ Cu $>$ Cr $>$ Pb, consistent with crustal
material, construction and resuspended road dust; enrichment factors flagged Cu, Zn and Cr
at traffic- and industry-influenced sites; PCA separated construction/road-dust, vehicular,
biomass-or-food and localised industrial factors. Hazard quotients ordered
Cr $>$ Pb $>$ Cu $>$ Zn with hazard indices below 1 for adults and children. The method is
the contribution: leaf-deposited dust is a zero-power, zero-maintenance passive sampler
that resolves metal-enriched microenvironments a sparse reference network cannot. Its
limits are equally clear --- deposition velocity varies with leaf morphology, species and
washoff history so concentrations are not convertible to airborne mass, integration time is
uncontrolled, and an \textbf{inhalation hazard quotient computed from deposited dust
requires an airborne concentration that is never observed here}.}
{10.1007/s10661-026-15906-w}

\paperentry{Sen et al.\ 2026}
{Genomic and phenotypic classification of Staphylococcus species isolated from Delhi air}
{Environ Microbiol Rep}
{Outdoor culturable bioaerosol sampling at six Delhi sites in November 2019 --- the peak
pollution month, pre-pandemic --- using a Merck MAS Eco-100. \textit{S.\ gallinarum},
\textit{S.\ haemolyticus} and \textit{S.\ arlettae} were recovered and characterised by
whole-genome sequencing, phenotypic microarray, antibiotic-resistance profiling and biofilm
formation in the presence of black carbon. Comparative genomics found the urease cluster
\textit{ureABCDEFG} in \textit{S.\ arlettae} Naraina and \textit{S.\ haemolyticus} Yojna
Vihar 1, unusual because urease is normally reported negative in both species and may
confer a niche fitness advantage. The black-carbon biofilm assay is the link to PM and the
part worth following, but \textbf{no quantitative result for it is reported}. Culture-based
sampling recovers a small biased fraction of the airborne microbiome, six sites in one
month give no seasonal or spatial contrast, and no particle mass was measured alongside
the biology, so the association with the pollution season is contextual.}
{10.1111/1758-2229.70411}


\band{Burden, policy \& mitigation\hfill 3 records}{Coral}

\paperentry{Ju et al.\ 2026}
{Premature mortality and economic losses attributable to landscape fire-related PM$_{2.5}$ in ten countries and territories}
{Lancet Planet Health}
{1{,}366 communities, ten countries and territories, \textbf{47{,}844{,}851 deaths},
2000--19, with community-specific associations estimated by instrumental-variable models
and pooled by random-effects meta-analysis. Each 1\,$\mu$g/m$^3$ of 2-day mean
landscape-fire PM$_{2.5}$ raised mortality \textbf{0.60\% (95\% CI 0.19--1.00)} and
mortality-related economic losses \textbf{0.59\% (0.21--0.97)}, yielding
\textbf{906{,}869 deaths and US\$994.5 billion} attributable (attributable fractions 1.90\%
and 1.64\%). The disproportion is the headline: fire PM$_{2.5}$ is \textbf{19.0\% of
ambient mass but 30.6\% of PM$_{2.5}$-attributable deaths}. IV estimates exceeded
conventional association estimates by a \textbf{median 56.0\% (2.2--111.3)}, so classical
designs are biased toward the null here. The weak point is that the instrument is not
named in the abstract and exclusion-restriction validity --- against co-emitted gases and
weather --- carries the causal claim entirely; value-of-statistical-life transfer across
ten economies adds uncertainty that is not propagated into the intervals.}
{10.1016/j.lanplh.2026.101506}

\paperentry{Bao et al.\ 2026}
{Distribution-sensitive assessment of air-quality and health benefits from clean-heating pilots in northern China}
{J Environ Manage}
{88 clean-heating pilot cities, 2013--23, evaluated with machine-learning weather
normalisation and the augmented synthetic control method against city-specific no-policy
counterfactuals. The distributional result is the point: the policy cut the heating-season
\textbf{mean only 3.01\,$\mu$g/m$^3$} but \textbf{P95 by 5.76, P99 by 5.96 and exceedance
days by 7.92} per season --- the tail moved roughly twice as far as the mean, which a
mean-only evaluation misses entirely. Health impact assessment gave \textbf{12{,}454 fewer
attributable all-cause deaths (95\% CI 8{,}564--16{,}222)}, with 88.9\% from the 2017--18
rollout cohorts and 56.2\% from the top decile of cities; high-pollution days were 27.8\%
of improvement days but \textbf{48.8\% of the mortality difference}. Benefit is therefore
concentrated in both time and space. ASCM counterfactuals depend on donor-pool
comparability during a decade of nationwide decline, and the mortality figure is a
transfer through an exposure-response function, not an observed death count.}
{10.1016/j.jenvman.2026.130960}

\paperentry{Mairghani et al.\ 2026}
{An umbrella review of public health interventions to mitigate air pollution in high-income countries}
{Rev Environ Health}
{Eight systematic reviews spanning city and transport policy, household exposure reduction
and emission-reduction modelling. Transport interventions gave consistent cardiovascular
benefit and fewer traffic injuries; portable air cleaners cut indoor PM$_{2.5}$ with modest
blood-pressure improvement; modelling projected a \textbf{median $\sim$1.5\% mortality
reduction} under large-scale mitigation. The ordering is the useful output --- structural
policy beat individual and household measures, which were complementary rather than
substitutive. Two limits constrain its use. \textbf{Equity was rarely assessed and
inconsistent where it was}, so the review cannot say who benefits; and the 1.5\% figure is
modelled rather than measured, mixing evidence tiers inside one summary. Eight reviews, all
high-income, so none of it transfers to the settings carrying the largest burden.}
{10.1515/reveh-2026-0061}

\band{Respiratory \& allergic\hfill 2 records}{Sage}

\paperentry{Meng et al.\ 2026}
{Exposure-response thresholds for ambient PM$_{2.5}$ and SO$_2$ vs COPD and small airway dysfunction, Yunnan}
{Front Public Health}
{11{,}095 residents aged 20+, multistage random cluster sampling 2021--22, with spirometry
in questionnaire-screened high-risk individuals. COPD prevalence \textbf{27.2\%}, small
airway dysfunction \textbf{53.6\%}. Restricted cubic splines gave significant nonlinearity
for both pollutants with thresholds at \textbf{21.33\,$\mu$g/m$^3$ PM$_{2.5}$} and
8.67\,$\mu$g/m$^3$ SO$_2$; above the PM$_{2.5}$ threshold, COPD and SAD prevalence were
33.4\% and 58.4\%. Random-forest importance ranked air pollution above smoking and biomass
fuel --- a strong claim from a cross-sectional design that should be read as variable
importance, not effect size. No temporality, exposure assigned without residential history,
and the enriched high-risk sampling frame means these prevalences do not transfer to a
general population. A 21\,$\mu$g/m$^3$ threshold sits well above the WHO 5\,$\mu$g/m$^3$
guideline and close to the regional mean, so it is plausibly a data-support artefact
rather than a biological inflection.}
{10.3389/fpubh.2026.1874156}

\paperentry{Vance et al.\ 2026}
{Greenspace, air pollution and respiratory health outcomes: a systematic review of cohort studies}
{Scand J Public Health}
{14 cohort studies to May 2026, all rated good on the Newcastle-Ottawa Scale,
predominantly European. Greenspace was generally protective against air-pollution-associated
respiratory mortality, and the methodological finding is the transferable one: \textbf{effects
appeared more consistently through mediation than moderation}, implying greenspace acts by
improving local air quality rather than by raising individual resilience. That has a direct
modelling consequence --- if the pathway is mediation, greenspace and PM$_{2.5}$ should not
be entered as independent covariates without care. No pooled estimate is reported, so the
synthesis is narrative and magnitude is unavailable; greenspace and pollution are both
strongly confounded by socioeconomic position; and near-complete European coverage limits
transfer.}
{10.1177/14034948261481328}


\band{Occupational \& indoor\hfill 2 records}{Clay}

\paperentry{Tam et al.\ 2026}
{Toxic air contaminant emissions from residential and commercial fossil gas appliances: a systematic review}
{Curr Environ Health Rep}
{25 studies, the first systematic synthesis of emission factors across gas appliance types.
Nonzero emissions documented for \textbf{at least 41 toxic air contaminants} via both
leakage and combustion pathways. Only formaldehyde and benzene had enough data for
quantitative characterisation, and their emission factors \textbf{spanned two orders of
magnitude} across studies and operating conditions --- a spread that makes any
single-value inventory estimate indefensible. Coverage is badly skewed toward cooking
stoves; heaters, water heaters and dryers are barely studied despite longer duty cycles.
The operationally interesting claim, correctly stated as emerging rather than established,
is that \textbf{methane and CO measurements may proxy for leakage- and combustion-related
TAC emissions} --- which is exactly the kind of cheap-channel inference target a low-cost
indoor node can carry.}
{10.1007/s40572-026-00562-6}

\paperentry{Yoshizaki et al.\ 2026}
{Differences in the effectiveness of smoke evacuator types for reducing surgical smoke: a randomized controlled trial}
{J Am Coll Surg}
{Randomised, double-blind trial (UMIN000047033), 64 laparotomy patients at the University
of Tokyo Hospital, Feb 2022 to May 2023, allocated 1:1:1 to electrocautery-incorporated
evacuation at 8\,mm tip distance ($n=22$), 18\,mm ($n=23$) or hose-type ($n=21$). The 8\,mm
system cut particle counts \textbf{96--97\% across 0.3--5.0\,$\mu$m} against hose-type;
\textbf{the 18\,mm arm did not reach significance on particles}, so effectiveness collapses
over a 10\,mm change in capture distance --- the practically important finding. Median
acetaldehyde was 1.90, 2.70 and 7.90\,$\mu$g/m$^3$ in the 8\,mm, 18\,mm and hose arms, both
electrocautery arms below hose, with formaldehyde behaving similarly. Endpoints are area
concentrations rather than personal breathing-zone exposure and no health outcome was
measured; single centre, one procedure type. The distance sensitivity is the transferable
lesson for any point-source capture or near-field sensing geometry.}
{10.1097/xcs.0000000000001943}

\band{Mechanistic toxicology\hfill 2 records}{Violet}

\paperentry{Ma et al.\ 2026}
{Combined PM$_{2.5}$ and microplastics exposure exacerbates allergic asthma via OPA1/SLC7A11-mediated ferroptosis}
{Ecotoxicol Environ Saf}
{Ovalbumin-induced murine asthma plus OVA-stimulated MRE cells exposed to PM$_{2.5}$ and/or
microplastics. Co-exposure amplified inflammatory infiltration, goblet-cell proportion and
lung-function impairment beyond either agent alone, with raised ROS and MDA and reduced
T-AOC, CAT and SOD. Ferroptosis markers moved coherently --- Fe$^{2+}$ accumulation, lipid
peroxidation, NCOA4/FTH1/ACSL4 up, SLC7A11/GPX4/GCLC/GSS down --- alongside mitochondrial
depolarisation, ATP depletion and reduced respiratory complex activity. \textbf{OPA1
knockdown and ferrostatin-1 both attenuated the injury}, two independent interventions
converging on one pathway, which is the strongest part of the design. What it cannot
support: no dose-response against realistic ambient concentrations, and instillation or
high-dose culture routinely exceeds inhaled doses by orders of magnitude. PM$_{2.5}$
composition is unspecified, so this is one sample, not a class.}
{10.1016/j.ecoenv.2026.120808}

\paperentry{Gulyaev et al.\ 2026}
{Polyphenols and inhalation-induced oxidative stress in chronic lung diseases: a systematic review}
{Biomed Rep}
{Systematic review to January 2026 across PubMed, Cochrane and Web of Science, with RoB 2
and SYRCLE risk-of-bias assessment. Preclinical evidence was consistent --- attenuated
oxidative damage, suppressed inflammatory signalling and enhanced antioxidant defence,
largely via Nrf2 activation and NF-$\kappa$B modulation, with histopathological improvement
across several exposure models. \textbf{Clinical evidence was not}: some trials improved
oxidative and inflammatory markers or symptom scores, but spirometric effects were
inconsistent. That preclinical-to-clinical gap is the familiar signature of antioxidant
intervention research, and the authors say so. Heterogeneity precluded meta-analysis, so
there is no pooled estimate; trials were small and short. Recorded as the mitigation-side
counterpart to the exposure literature, not as evidence that dietary intervention
substitutes for reducing exposure.}
{10.3892/br.2026.2194}


% 2026-09-16 proof: this band header rendered alone at the foot of p8 with no entry
% under it. Structural break rather than shaving -- see PIPELINE 13/14/15 Sep.
\clearpage

\band{Neuro / mental health\hfill 1 record}{Deep}

\paperentry{Zhang et al.\ 2026}
{Household solid fuel use, kitchen ventilation and depressive symptoms among older adults in China}
{Front Public Health}
{11{,}697 CLHLS participants aged 65+, 2018 wave. Depressive symptoms (CES-D-10) in
\textbf{26.87\%}, solid fuel use in 29.67\%. Solid fuel carried \textbf{OR 1.25 (95\% CI
1.09--1.44)} against clean fuel and \textbf{B = 0.57 (0.33--0.82)} on the continuous score
--- consistent but small, roughly a fifth of a scale point per unit. The charcoal subgroup
showed a larger estimate on $n=39$ exposed and the authors correctly flag it as unreliable.
\textbf{No interaction was found between fuel type and kitchen ventilation}, which
undercuts the paper's own framing that ventilation modifies the exposure. Exposure is
entirely self-reported category with no PM$_{2.5}$ measurement, so dose is unknown and
misclassification is certain to be substantial; and the cross-sectional design cannot
exclude reverse causation, since depressed older adults may differ in both household energy
choice and reporting.}
{10.3389/fpubh.2026.1920233}

\band{Cardiovascular \& metabolic\hfill 1 record}{Teal}

\paperentry{Shahriary et al.\ 2026}
{Cardiovascular and pregnancy-related health effects of ultrafine particulate matter}
{Toxicol Lett}
{Mechanistic synthesis of experimental, controlled-human-exposure and epidemiological
evidence linking ultrafine PM to cardiovascular injury, with emphasis on hypertensive
disorders of pregnancy. Pathways covered: pulmonary and systemic inflammation, oxidative
stress, mitochondrial dysfunction, autonomic imbalance, endothelial injury and disturbed
angiogenic balance, extended to maternal cardiovascular adaptation and placental vascular
development. The review's honest core is its attribution problem --- it states plainly that
\textbf{UFPM-specific findings are hard to separate from PM$_{2.5}$, diesel exhaust
particles, engineered nanoparticles and traffic-related pollution generally}, and that
direct developmental evidence is largely animal. Uncertainties named: particle
translocation, exposure metric, critical gestational windows. No pooled estimates. The
exposure-metric point is the operational one --- UFP is a number-concentration phenomenon
and mass-based networks, including nearly all low-cost optical nodes, do not measure it.}
{10.1016/j.toxlet.2026.113203}

\band{Reproductive \& developmental\hfill 1 record}{Sky}

\paperentry{Ammons et al.\ 2026}
{Environmental chemical exposures and ovarian and endometrial cancers: a systematic review}
{Environ Res}
{145 studies screened --- 54\% ovarian-cancer only, 53\% North American, 40\% cohort ---
with outdoor air pollution the single most studied exposure class (27 studies, 18\%).
The pooled PM$_{2.5}$ result is a clean null for ovarian cancer: \textbf{HR 1.02 (95\% CI
0.91--1.13)} per 5\,$\mu$g/m$^3$ across four studies, and \textbf{0.94 (0.81--1.09)}
restricted to the three with co-pollutant adjustment. Endometrial cancer was weakly
positive and implausibly precise --- 1.02 (0.99--1.05) overall but \textbf{1.01
(1.00--1.01)} with co-pollutant adjustment and 1.01 (1.01--1.02) for incident cases only.
\textbf{A confidence interval that narrow from five studies signals near-zero
between-study heterogeneity or shared cohorts rather than genuine precision}, and no
$I^2$ is reported. Worth recording as a negative: gynaecologic cancer is one endpoint where
the PM$_{2.5}$ signal does not appear, against a literature that increasingly reports
positive associations wherever it looks.}
{10.1016/j.envres.2026.125661}

\clearpage

% ==================================================================== TABLE
\section*{\sffamily\bfseries\color{Deep}Corpus register}
\vspace{-1.5mm}{\color{Amber}\rule{\linewidth}{1.1pt}}\vspace{2.5mm}

{\fontsize{6.7}{8.4}\selectfont
\setlength{\tabcolsep}{2.4pt}
\begin{longtable}{@{}>{\RaggedRight}p{25mm}>{\RaggedRight}p{41mm}>{\RaggedRight}p{27mm}%
>{\RaggedRight}p{23mm}>{\RaggedRight}p{18mm}>{\RaggedRight}p{15mm}@{}}
\toprule
\textsf{\textbf{First author}} & \textsf{\textbf{Focus}} & \textsf{\textbf{Journal}} &
\textsf{\textbf{Design}} & \textsf{\textbf{Metric}} & \textsf{\textbf{Region}}\\
\midrule
\endfirsthead
\toprule
\textsf{\textbf{First author}} & \textsf{\textbf{Focus}} & \textsf{\textbf{Journal}} &
\textsf{\textbf{Design}} & \textsf{\textbf{Metric}} & \textsf{\textbf{Region}}\\
\midrule
\endhead
\multicolumn{6}{@{}l@{}}{\cellcolor{Mist}\textsf{\textbf{\textcolor{Deep}{Neuro / mental health}}}}\\[0.6mm]
\href{https://doi.org/10.1001/jamanetworkopen.2026.25015}{Dales} & Wildfire-sourced PM$_{2.5}$ and ED visits for migraine and other primary headache & \textit{JAMA Netw Open} & Case-crossover & PM$_{2.5}$ (wildfire) & Canada\\
\href{https://doi.org/10.1093/ageing/afag215}{Glover} & Umbrella meta-review of urban environmental exposures and cognitive health & \textit{Age Ageing} & Umbrella review / meta-analysis & PM$_{2.5}$ & Global\\
\href{https://doi.org/10.1016/j.nsa.2026.107019}{Goossens} & In utero air pollution exposure and pre-/postnatal brain development & \textit{Neurosci Appl} & Narrative review & PM$_{2.5}$ / PM$_{10}$ / PM$_{\mathrm{coarse}}$ & Global\\
\href{https://doi.org/10.1016/j.isci.2026.116617}{Ma} & PM$_{2.5}$-induced neurotoxicity and depression: exposure risks and mechanisms & \textit{iScience} & Narrative review & PM$_{2.5}$ + components & Global\\
\href{https://doi.org/10.3760/cma.j.cn112150-20250804-00752}{Pan} & Maternal PM exposure and offspring neurodevelopment: effects and mechanisms & \textit{Zhonghua Yu Fang Yi Xue Za Zhi} & Narrative review & PM (bulk) & China\\
\addlinespace[0.9mm]
\multicolumn{6}{@{}l@{}}{\cellcolor{Mist}\textsf{\textbf{\textcolor{Teal}{Cardiovascular \& metabolic}}}}\\[0.6mm]
\href{https://doi.org/10.3760/cma.j.cn112338-20251120-00843}{Huang} & PM$_{2.5}$ and dyslipidaemia in rural Chinese children; dietary effect modification & \textit{Zhonghua Liu Xing Bing Xue Za Zhi} & Cross-sectional (multistage) & PM$_{2.5}$ & China\\
\href{https://doi.org/10.32394/pe/218966}{Kazmierska} & Air pollution and cardiovascular disease: review of current data & \textit{Przegl Epidemiol} & Narrative review & PM$_{2.5}$ / PM$_{10}$ & Global\\
\href{https://doi.org/10.1007/s00125-026-06814-2}{Xi} & Air pollution, built environment and temperature vs. prediabetes phenotypes & \textit{Diabetologia} & Prospective cohort & PM$_{2.5}$ & Germany\\
\href{https://doi.org/10.1161/JAHA.125.046917}{Yang} & Short-term PM$_{2.5}$ and hospitalisation/ED use in atrial fibrillation & \textit{J Am Heart Assoc} & Case-crossover & PM$_{2.5}$ (wildfire-influenced) & USA\\
\addlinespace[0.9mm]
\multicolumn{6}{@{}l@{}}{\cellcolor{Mist}\textsf{\textbf{\textcolor{Sky}{Reproductive \& developmental}}}}\\[0.6mm]
\href{https://doi.org/10.1016/j.jes.2025.10.043}{Chen} & Prenatal air pollution, fetal growth trajectory phenotypes and child neurodevelopment & \textit{J Environ Sci (China)} & Prospective cohort & PM$_{2.5}$ / PM$_{10}$ & China\\
\href{https://doi.org/10.1016/j.ecoenv.2026.120566}{Huang} & Joint prenatal bisphenol + air-pollutant mixture exposure and preschool IQ & \textit{Ecotoxicol Environ Saf} & Prospective cohort (mixtures) & PM$_{2.5}$ / PM$_{10}$ & China\\
\href{https://doi.org/10.1007/s40572-026-00552-8}{Peterson} & Ambient air pollution, PAHs and pubertal development: critical review & \textit{Curr Environ Health Rep} & Critical review & PM$_{2.5}$ components & Global\\
\href{https://doi.org/10.1016/j.envres.2026.125313}{Xie} & Long-term PM$_{2.5}$ and timing of pubertal development in boys and girls & \textit{Environ Res} & Prospective cohort & PM$_{2.5}$ & China\\
\addlinespace[0.9mm]
\multicolumn{6}{@{}l@{}}{\cellcolor{Mist}\textsf{\textbf{\textcolor{Sage}{Respiratory \& allergic}}}}\\[0.6mm]
\href{https://doi.org/10.1007/s00484-026-03289-z}{Bayar Muluk} & Climate change and otolaryngology: airway disease impacts and emission reduction & \textit{Int J Biometeorol} & Narrative review & PM$_{2.5}$ & Global\\
\href{https://doi.org/10.4103/mgr.MEDGASRES-D-25-00164}{Liao} & Short-term air pollution, respiratory symptoms and lung function in asthma & \textit{Med Gas Res} & Panel study & PM$_{2.5}$ / PM$_{10}$ & China\\
\href{https://doi.org/10.3390/epidemiologia7040097}{Sanchez de Prada} & Ambient PM and asthma-coded admissions in children 0-3 y, Spain & \textit{Epidemiologia} & Ecological panel & PM$_{2.5}$ / PM$_{10}$ & Spain\\
\href{https://doi.org/10.1016/j.jhazmat.2026.143076}{Yang} & Allergenic vs. toxic risk of airborne seed hairs compared with ambient PM$_{2.5}$ & \textit{J Hazard Mater} & Field measurement & PM$_{2.5}$ + bioaerosol & China\\
\addlinespace[0.9mm]
\multicolumn{6}{@{}l@{}}{\cellcolor{Mist}\textsf{\textbf{\textcolor{Violet}{Mechanistic toxicology}}}}\\[0.6mm]
\href{https://doi.org/10.1016/j.jes.2025.12.074}{Ding} & PM$_{2.5}$ promotes vascular calcification via miR-27b-5p/HES1 signalling & \textit{J Environ Sci (China)} & In vitro + organoid + murine & PM$_{2.5}$ & China\\
\href{https://doi.org/10.1007/s00360-026-01698-5}{Maraschi} & Osmoregulatory disruption and metal bioaccumulation from settleable atmospheric PM & \textit{J Comp Physiol B} & Controlled exposure (ecotox) & Settleable APM & Brazil\\
\href{https://doi.org/10.1016/j.jes.2025.11.018}{Wu} & Macrophage NLRP3 inflammasome in real-ambient PM-induced hepatic glucose dysregulation & \textit{J Environ Sci (China)} & Real-ambient murine + co-culture & Real-ambient PM & China\\
\addlinespace[0.9mm]
\multicolumn{6}{@{}l@{}}{\cellcolor{Mist}\textsf{\textbf{\textcolor{Clay}{Occupational \& indoor}}}}\\[0.6mm]
\href{https://doi.org/10.3389/fmedt.2026.1796828}{Chen} & Nested suction instrument to cut intraoperative surgical-smoke PM$_{2.5}$ & \textit{Front Med Technol} & Bench / device experiment & PM$_{2.5}$ (surgical smoke) & China\\
\href{https://doi.org/10.1093/annweh/wxag055}{Cole} & Determinants of respirable crystalline silica exposure in Australian tunnelling & \textit{Ann Work Expo Health} & Observational exposure study & Respirable dust / RCS & Australia\\
\href{https://doi.org/10.1016/j.jhazmat.2026.143081}{Ding} & Stage-resolved hair Zn-isotope dynamics as a firefighting combustion-exposure tracer & \textit{J Hazard Mater} & Repeated-measures biomarker & Combustion-derived PM & China\\
\href{https://doi.org/10.3390/jox16040119}{Ramsden} & Xenobiotic hazards in aircraft cabin air & \textit{J Xenobiot} & Review + risk economics & Ultrafine particles & Global\\
\addlinespace[0.9mm]
\multicolumn{6}{@{}l@{}}{\cellcolor{Mist}\textsf{\textbf{\textcolor{Amber}{Sensing, forecasting \& instrumentation}}}}\\[0.6mm]
\href{https://doi.org/10.1016/j.jenvman.2026.130515}{Shi} & Adaptive spatiotemporal graph transformer for multi-site multi-step PM$_{2.5}$ forecasting & \textit{J Environ Manage} & Deep-learning model & PM$_{2.5}$ & China / India\\
\href{https://doi.org/10.3390/bios16070380}{Tsuchiya} & Seasonal and indoor-outdoor characteristics of fluorescent aerosol particles in offices & \textit{Biosensors} & Real-time sensor campaign & Bioaerosol / coarse & Japan\\
\addlinespace[0.9mm]
\multicolumn{6}{@{}l@{}}{\cellcolor{Mist}\textsf{\textbf{\textcolor{Amber}{Exposure assessment \& modelling}}}}\\[0.6mm]
\href{https://doi.org/10.1029/2025GH001621}{Solorzano-Araque} & Multiscale exposure modelling: LOTOS-EUROS + mobility surveys + particle toxicity & \textit{GeoHealth} & CTM + mobility integration & PM$_{2.5}$ / PM$_{10}$ & Colombia\\
\addlinespace[0.9mm]
\multicolumn{6}{@{}l@{}}{\cellcolor{Mist}\textsf{\textbf{\textcolor{Coral}{Burden, policy \& mitigation}}}}\\[0.6mm]
\href{https://doi.org/10.1002/hsr2.72886}{Karamzad} & Ischaemic heart disease burden in MENA 1990-2021; PM attributable fraction 28.2\% & \textit{Health Sci Rep} & GBD secondary analysis & Ambient PM & MENA\\
\href{https://doi.org/10.1371/journal.pone.0347989}{Rees} & HEPA filtration to reduce respiratory infection in care homes & \textit{PLoS One} & RCT process evaluation & Filtered PM & UK\\
\href{https://doi.org/10.1016/j.jes.2026.01.009}{Shan} & Policy-driven emission reductions and post-pandemic rebound, high-resolution inventories & \textit{J Environ Sci (China)} & Emission inventory & PM emissions & China\\
\addlinespace[0.9mm]
\multicolumn{6}{@{}l@{}}{\cellcolor{Mist}\textsf{\textbf{\textcolor{Slate}{Other clinical endpoints}}}}\\[0.6mm]
\href{https://doi.org/10.3390/dj14070440}{Lupita} & PM$_{2.5}$, preventive-care capacity and caries/edentulism burden across the EU & \textit{Dent J} & Ecological panel & PM$_{2.5}$ & EU-27\\
\href{https://doi.org/10.1016/j.jdin.2026.06.008}{Park} & Long-term PM exposure and risk of androgenetic alopecia & \textit{JAAD Int} & Nationwide cohort & PM$_{2.5}$ / PM$_{10}$ & South Korea\\
\href{https://doi.org/10.3171/2026.3.PEDS2615}{Sun} & Environmental exposure and area deprivation vs. paediatric hydrocephalus outcomes & \textit{J Neurosurg Pediatr} & Registry cohort + geospatial & PM$_{2.5}$ (satellite-derived) & USA\\
\addlinespace[0.9mm]
\bottomrule
\end{longtable}}


\vspace{4mm}
\begin{keybox}
{\sffamily\bfseries\footnotesize\color{Deep}Provenance and method}\\[1.2mm]
{\fontsize{7.2}{8.8}\selectfont
Entry window \textbf{2026-09-16 $\to$ 2026-09-16}, contiguous with the 15 September issue;
no gap. Built on 17 September under the standing rule that the current day is not opened
before 22:00 local.
\textbf{Sources.} PubMed \texttt{[Date - Entry]} run as two separate queries --- health
axis and instrumentation axis --- through both the local harvester and the connector:
\textbf{16 unique candidates}, the harvester a strict superset of the connector this time
(16 vs 14), reversing the 15 September finding that neither client alone was sufficient.
Europe PMC \texttt{CREATION\_DATE} 56. Crossref windowed by ISSN across the tracked
non-PubMed journals \textbf{66 deposits}, which carried all four remote-sensing records.
\textbf{arXiv returned HTTP 406 again}; three header variants were tested this run
(HTTPS with an Atom \texttt{Accept}, a browser user-agent, and plain HTTP) and all three
failed identically, so the cause is not header negotiation and the leg should now be
routed through a connector rather than repaired in \texttt{harvest.py}. OpenAlex HTTP 429,
as on every recent run.
\textbf{Consensus returned nothing admissible.} Two sweeps, 20 records: \textbf{seven were
already in the archive} on a title check, and \textbf{all five that survived title
screening and were then date-checked failed} --- \texttt{10.3390/s26010280} created
2026-01-02, \texttt{10.1021/acs.est.6c00095} 2026-05-08,
\texttt{10.1016/j.scitotenv.2026.181915} 2026-05-27, \texttt{10.1002/met.70228}
2026-08-22, and the India AOD fusion work a December 2025 egusphere preprint. Consensus
has no entry-date filter, so every result needs an independent Crossref \texttt{created}
check; on this run the admitted yield was zero for the first time.
\textbf{Europe PMC AGRICOLA defect.} Two AGR records returned inside a 16 September
\texttt{CREATION\_DATE} window resolve to \textbf{pubYear 2021 and 2014} --- the Fishlake
UAS emission-sampling paper and the Hodas outdoor-to-indoor transport model. The AGR leg
retro-indexes, so a publication-year check is now mandatory before any AGR record is
summarised. Both were rejected.
\textbf{Screening.} 138 raw records across the four legs, 130 unique after cross-source
dedupe, of which 8 were caught by the DOI gate against \texttt{seen.json}.
\textbf{21 shipped, 27 excluded, 3 held.} The three held are same-day Elsevier deposits
--- Moscow long-term black carbon, the APR hybrid-imputation multi-target random forest,
and the \textit{Building and Environment} indoor CO$_2$ sensor-siting study --- whose DOIs
resolve but for which no abstract exists in Crossref, Europe PMC or Semantic Scholar. All
three are on-axis and are deliberately kept out of \texttt{seen.json} so a later issue can
carry them; this is the second consecutive issue in which the same Elsevier
no-abstract-on-deposit pattern has cost tier-A candidates.
\textbf{Trial watch.} Two registry records updated in window. NCT07540000 (INTERLUNG,
Johns Hopkins, Bhaktapur, Nepal, PHASE4, $n=800$) was already stored; endpoints re-checked
and unchanged --- pre-bronchodilator FEV$_1$ z-score at 40 months powered for 0.25\,SD,
plus SGRQ, the strongest genuinely respiratory outcome set in the accumulated file.
NCT07793123 (Robertet SA, oral antioxidant vs placebo, skin radiance, $n=60$) is new and
kept, but is air-pollution-adjacent only: \textbf{all eight endpoints are dermatological or
biophysical and no PM variable is measured} --- the third consecutive topical hit whose PM
framing no measured variable carries. \texttt{trials.json} was also \textbf{reconciled}:
the running \texttt{trials} tally of 28 flagged on 15 September has been replaced by the
derived count of persisted NCT records (12), and the counter must now be derived, never
incremented.
\textbf{Labels.} Subtopic, design, particle-metric and geography are assigned from title,
abstract, keywords and MeSH; single-label by dominant contribution, a judgement rather than
a controlled vocabulary. Five new design labels and one new geography key were registered
in \texttt{plots.py} in the same edit that wrote them onto records.
Effect estimates are transcribed verbatim from abstracts; the two landscape-fire percent
changes are converted to risk-ratio equivalents and flagged as such in the forest-plot
caption, and nothing else is rescaled.
Content derived from PubMed and Crossref metadata; abstracts remain the copyright of their
respective publishers.\par}
\end{keybox}

\end{document}
